\subsection{Games and logical interpretations}
\label{sec:interpretations}

\subsubsection{Full strategies and executable plans}
\label{sec:policies}

A full strategy records behavior at every well-typed history, including
histories ruled out by the player's own earlier choices. We retain this
object for strategic analysis and use the endpoint plans of \eqref{dep:eq:4} for
execution. The distinction is the standard one between strategies and
plans of action \cite{Rubinstein1991}. It is not a choice of solution concept.

Define the full local strategy type by
\[
 \mathsf S_i(\End\langle R_j\rangle_j)=R_i,\qquad
 \mathsf S_i(\Q_r(x:A).H)=
 \begin{cases}
 A\times\displaystyle\prod_{x:A}\mathsf S_i(H),&i=r,\\[3pt]
 \displaystyle\prod_{x:A}\mathsf S_i(H),&i\ne r.
 \end{cases}
\]
At an owned choice, $(a,f)$ chooses $a$ and retains $f(x)$ for every
alternative $x$, including those not chosen. Escard\'o and Oliva use this
move-and-all-continuations construction in their dependent strategy type
\cite[Definition~4.1]{EscardoOliva2023}. Our per-role version distributes
those decisions among the owners.

Reduction to an executable plan is recursive:
\[
 \mathsf{red}_i(a,f)
   =(a,\mathsf{red}_i(f(a)))\quad(i=r),\qquad
 \mathsf{red}_i(f)=\lambda x.\mathsf{red}_i(f(x))\quad(i\ne r).
\]
At a terminal it is the identity on $R_i$.
Write $\Run_G^{\mathsf{full}}$ for full-strategy execution: it uses $a$
from the owner's pair and continues with
$f(a)$ there and each receiver's application at $a$. Terminal returns are
collected as before.

This construction has both a syntactic interpretation in $\D$ and a
set-theoretic interpretation over a model of the public data. In the latter,
all set-theoretic dependent functions are admitted; they need not be
definable data programs.
The finite rational analysis in \cref{sec:rational} uses these complete
set-theoretic strategy spaces, whereas executable plans in $\D$ provide
the program interface for the public execution core.

\begin{lemma}[Execution adequacy for strategy reduction]
\label{thm:adequacy}
For every full profile $\sigma$,
$\Run_G^{\mathsf{full}}(\sigma)
 =\Run_G((\mathsf{red}_i(\sigma_i))_i)$:
execution gives the same trace and return record as its reduced endpoint
plans.
\end{lemma}
\begin{proof}
At a terminal both executions use the same returns. At a choice both use
the same owner's $a$. Reducing each full continuation at $a$ gives exactly
the residual reduced plans. Induct on the common continuation.
\end{proof}

Equality of this execution is weaker than equality of strategic behavior.
For example, take two Boolean choices both owned by Alice. Alice wins iff
$x=0$ or $y=1$; Bob has the complementary objective and owns no move.
Two full Alice strategies choose $x=0$ and $y=0$ there, but prescribe,
respectively, $y=0$ and $y=1$ after the alternative $x=1$.
They reduce to the same plan and execute the same trace $(0,0)$.
Only the second is subgame perfect: in the first, Alice would improve
from losing to winning by changing $y$ at the unreached subgame $x=1$.
Both are Nash equilibria, since the actual payoff is already maximal.

Thus subgame perfection does not factor through $\mathsf{red}$.
Retaining full strategies avoids that loss without changing the smaller
execution interface.

\subsubsection{Objectives observe executions}

An objective can be any specified predicate $W_i(e)$ on $e:\Exec(G)$; a utility
can be a function $u_i(e)$ into an ordered payoff domain. An outcome tuple
$\langle[W_i(e)]\rangle_i$ is therefore one readout of the complete execution.
No exclusivity or agreement condition is imposed on the $W_i$.

For example, after \eqref{dep:eq:1}, Alice may want the selected entry to be true, Bob may
want it to be false, and Carol may want the first position. These predicates
can conflict while the interaction remains well formed and executable.
Taking $R_i=\Unit$ does not certify any of them. Replacing $R_i$ by a
particular evidence type makes that proof a requirement of endpoint $i$;
it changes the interface explicitly.

Objectives are predicates on executions. Solution concepts are predicates
on strategies in a game. For a full profile $\sigma$, Nash equilibrium
requires, for every role $i$ and full unilateral deviation $\tau_i$,
\[
 u_i(\Run^{\mathsf{full}}(\sigma))\ge
 u_i(\Run^{\mathsf{full}}(\sigma[i:=\tau_i])).
\]
Subgame perfection requires the corresponding Nash condition in every
subgame. Neither condition is a predicate of the actual execution alone.

For example, Alice chooses $L$ or $R$, then Bob chooses $U$ or $D$.
Every $L$-trace pays $(1,1)$; $RU$ pays $(0,0)$ and $RD$ pays $(2,1)$.
Fix the actual choices $L,U$. If Bob prescribes $U$ after $R$, this is a
Nash profile. If he instead prescribes $D$ there, Alice can profitably
deviate to $R$. The actual trace and payoff are unchanged. Analysis must
retain the deviation structure as well as execution.

The interaction object therefore consists of the script, full strategy
spaces, and evaluation. Objectives and solution concepts are added at their
respective types. No guaranteeing or equilibrium condition defines the
formation of a role-owned choice.

\subsubsection{Classical quantifiers as a two-role interpretation}

Hintikka's evaluation game of first-order logic is recovered by fixing
a viewpoint and complementary objectives \cite{Hintikka1973}. We state this standard
specialization over finite nonempty structures, so it involves no choice of
noncomputable strategies over infinite domains.

Fix distinct roles $i,j$ and a prenex first-order sentence
$Q_1x_1\cdots Q_nx_n.\theta$, with quantifier-free matrix $\theta$.
Translate each existential binder to $\Q_i$, each universal binder to
$\Q_j$, and let all terminal result types be $\Unit$. Give role $i$ objective
$\theta$ and role $j$ objective $\neg\theta$. The positive objective belongs
to the selected role $i$, whichever role name is chosen.

\begin{proposition}[Classical evaluation-game correspondence]
In a finite nonempty first-order structure, the sentence is true iff role
$i$ has a policy that achieves $\theta$ against every policy of $j$.
\end{proposition}
\begin{proof}
Induct on the quantifier prefix. At an existential binder, truth requires
some value with a true continuation, and role $i$ can select that value.
At a universal binder, every value must have a true continuation; role $i$
can use the corresponding contingent continuation after observing $j$'s
choice. Finite branching permits all these choices to be assembled into
one policy. Conversely, a guaranteeing policy gives the existential witness
at an owned binder and a guaranteeing continuation after every opposing
choice. At the terminal the statement is satisfaction of $\theta$.
\end{proof}

This interpretation explains the existential/universal pair without assigning
permanent logical polarities to the roles of the global language. In a
general $K$-role interaction there is no distinguished opposing role.

\begin{proposition}[Role permutation is not general Boolean complement]
For $K>2$, no permutation of individual roles maps every one-hot outcome to
its coordinatewise Boolean complement.
\end{proposition}
\begin{proof}
A permutation preserves the number of nonzero coordinates. A one-hot vector
has one, while its complement has $K-1$.
\end{proof}

At $K=2$ these counts coincide, and swapping the complementary evaluation-game
roles has the familiar negation reading. For larger role sets, permutation
remains a symmetry of interaction. Logical complement of an objective remains
an operation in the objective language. The two operations have different
types and different laws.
